\begin{abstract}
Transistor network synthesis—finding a compact CMOS switching network that
correctly implements a given Boolean function—is a central step in
custom-cell design.
Existing exact and SAT-based methods are powerful for small functions but scale
poorly as the number of variables grows and cannot generalize across functions
without re-solving from scratch.
We present \method, a learning-based framework that casts transistor network
synthesis as a sequential graph-generation Markov decision process (MDP) and
trains a \sgpn (\SGPN) to solve it.
The policy operates over a unified graph comprising both the partial synthesized
network and a fixed SOP-based compensate network that encodes the target Boolean
function, allowing the GNN backbone to reason jointly about coverage obligations
and existing topology at every step.
Training proceeds in three phases:
Phase~1 performs supervised NLL pretraining on a dataset of known-optimal
transistor networks to initialize the policy;
Phase~2 applies PPO-based reinforcement learning with a GRPO group baseline and
an adaptive self-tightening reward to push the policy toward minimum transistor
count;
Phase~3 extends RL to joint pull-down/pull-up network (PDN/PUN) synthesis
with a physical placement reward derived from ASTRAN, a standard-cell layout
tool, to optimize cell width (CPP) directly.
On a benchmark of \TODO{N} Boolean functions with up to three variables,
\method{} achieves \TODO{X\%} synthesis success rate and matches or improves upon
the reference transistor count on \TODO{Y\%} of functions;
physical-aware finetuning further reduces the ASTRAN placement objective
by \TODO{Z\%} compared to the Phase~2 baseline,
demonstrating that a single policy can jointly optimize logical correctness,
transistor count, and physical layout quality.
\end{abstract}
